\section{Introduction}

For a long-lived deployed AI service, repeated inference can dominate its
operational energy, and the cost of each generated token depends strongly on
the request profile and serving configuration
\cite{samsi2023words,niu2025engines,fernandez2025energy}.  Batch and token
limits, chunked prefill, prefix caching, precision, and power policy
jointly determine throughput, time to first token (TTFT), time per output token
(TPOT), and energy.  Long prompts emphasize prefill compute, long responses
emphasize memory-intensive decode, and request concurrency changes batch
occupancy and queueing.  Choosing the serving configuration is therefore an
energy decision as well as a performance decision.

The obvious solution, serving the workload under every candidate and measuring
the result, is often unavailable.  H100/H200/B200 allocations are scarce,
a dense sweep spends the most valuable evidence on configurations that a CPU
model or short probe could have eliminated.  Yet cheap evidence is not hardware
truth.  A CPU model can mispredict a latency boundary even when its energy
ranking is useful, while a short run may not reproduce the batching and
queueing behavior of the full target trace.  The central systems question is
thus not merely
\emph{which configuration is best?}, but \emph{which
configuration--fidelity pair should be evaluated next under a GPU-hour budget?}

The measurement path itself is increasingly mature.  \bench{} provides
declarative, phase-aligned GPU, node, and system telemetry for inference
\cite{niu2026tokenpowerbench}.  BEAM jointly controls
batching and frequency using an offline energy--performance table, BOute applies
multi-objective Bayesian optimization to deployment choices, and OptiKIT uses
full hardware trials for serving tuning
\cite{lee2026beam,jiang2026boute,santavas2026optikit}.  Charon and a Meta
deployment study make broad search practical through calibrated simulation,
while PROMPTS adds agent reasoning to distributed configuration recommendation
\cite{yang2026charon,cho2026deployment,ding2026prompts}.  These systems optimize
a configuration, runtime action, or agent proposal.  They do not jointly choose
the cost and scale of the next observation while requiring measured evidence
before releasing an energy--SLO recommendation.

We present \system{}, an evidence-gated, multi-fidelity search framework built
on \bench{}.  A scenario specifies the model, workload, target GPU, serving
runtime, SLOs, candidate knobs, and measurement budget.  \system{} compiles a
typed feasible pool and selects actions $a=(x,e)$, where $x$ is a serving
configuration and $e$ is \textsc{Sandbox}, \textsc{Probe}, or
\textsc{Verify}.  These stages correspond to the artifact identifiers L0, L1,
and L4; the semantic names describe the three-stage path exercised here.  Its
\emph{Intent-Conditioned Pareto Information Gain} (\ipig{}) interface scores the
expected reduction in uncertainty about the verified SLO-feasible frontier per
GPU-hour.  An LLM maps natural-language intent to a typed subgoal and explains
failures; deterministic tools compile candidates, score actions, account for
cost, ingest telemetry, compute Pareto sets, and enforce final verification.
The LLM never generates energy estimates or receives release authority.

We evaluate both the evidence substrate and the fidelity ladder on one H100.
First, a frozen CPU-side model is tested on \BlindMeasurements{} blind
Qwen2.5-7B workload-transfer measurements.  Second, a preregistered
\ConfigurationCandidates{}-candidate serving grid varies sequence capacity,
batched-token limits, and chunked prefill.  Its \ConfigurationMeasuredRuns{}
balanced short-probe and target-workload measurements expose the exact boundary
the system is designed to manage: the Sandbox ranks energy well but fails for TTFT and
the Pareto decision, whereas the short Probe recovers the true verified frontier
with \LoneParetoRecall{} recall.  The full target run then removes two false
positives and verifies a
configuration that improves energy, throughput, TTFT, and TPOT over the expert
default.  In corrected empirical replay, \ipig{} improves oracle-hit rate and
GPU-hour cost over random and cost-blind acquisition, but not over a static
cheapest-first ladder.  The corpus is sealed before replay; its exhaustive cost
is ground-truth construction, not an agent-savings claim.

This paper makes four contributions:
\begin{itemize}
    \item We formulate energy-aware inference tuning as \emph{budgeted
    configuration--fidelity co-selection}: recover an SLO-feasible
    energy--performance frontier while minimizing real GPU-hours spent acquiring
    and verifying evidence.
    \item We design \sandbox{}, a three-stage evidence contract that combines
    CPU-first projection, a short measured probe, uncertainty, immutable
    provenance, and mandatory full-workload verification on one target GPU.
    \item We implement a bounded hybrid agent with an \ipig{} policy interface,
    typed executors, failure-aware reflection, hard budget and action guards,
    and a deterministic release gate.
    \item We provide \BlindMeasurements{} blind workload-transfer runs and a
    \ConfigurationMeasuredRuns{}-run configuration corpus.  The measurements
    quantify when the Sandbox is informative, when the agent must escalate to a
    short Probe, and how full verification converts a three-candidate
    provisional frontier into one verified configuration; sealed replay then
    compares four acquisition policies without spending additional GPU-hours.
\end{itemize}
